\documentclass[../../main.tex]{subfiles}
\begin{document}

{\bf [解答]}
$A$，$B$の$P$からの距離をそれぞれ$A(t)$，$B(t)$とすると，
\begin{align}
  A(t)=\int_0^t u\,dt+A(0)=ut \\
  B(t)=\int_0^t v\,dt+B(0)=\dfrac{1}{4}t^3-\dfrac{3}{2}t^2+25
\end{align}
である．$B(t)\le25\iff0\le t\le6$であるから，$A$が$B$に追いつくには，$y=A(t)$と$y=B(t)$が$0<t<6$で少なくとも一回交わればよい．$y$を消去して
\begin{align}
  &A(t)=B(t) \nonumber \\
  \iff &u=\dfrac{1}{4}t^2-\dfrac{3}{2}t+\dfrac{25}{t}\equiv f(t) \label{eq:1}
\end{align}
である．
\begin{align}
  f'(t)=\dfrac{1}{2}t-\dfrac{3}{2}-\dfrac{25}{t^2}=\dfrac{(t-5)(t^2+2t+10)}{2t^2}
\end{align}
から，\cref{tab:variation_ft}を得る．

\begin{table}[htb]
  \centering
  \caption{$f(t)$ の増減表}
  \label{tab:variation_ft}
  \begin{tabular}{c|c|c|c|c|c}
    $t$  & $0$ & $\dots$ & $5$ & $\dots$ & $6$ \\ \hline
    $f'$ &     & $-$     & $0$ & $+$     &     \\ \hline
    $f$  &     & $\searrow$ & $\dfrac{15}{4}$ & $\nearrow$ & $\dfrac{25}{6}$
  \end{tabular}
\end{table}

よって，$f(t)\to\infty\ (t\to+0)$と合わせて，グラフは\cref{fig:graph_ft}のようになる．

\begin{figure}[htb]
  \centering
  \begin{tikzpicture}[x=0.9cm, y=0.9cm, >=stealth,
    every node/.style={font=\small}]
    % 座標軸
    \draw[->, thick, gray!80!black] (-0.5, 0) -- (7.0, 0) node[right] {$t$};
    \draw[->, thick, gray!80!black] (0, -0.5) -- (0, 5.5) node[above] {$y$};
    \node[below left, font=\footnotesize] at (0, 0) {$\mathrm{O}$};

    % y = f(t) の曲線
    \draw[line width=1.2pt, blue!85!black, domain=1.8:6.2, samples=60]
    plot (\x, {0.25*\x*\x - 1.5*\x + 25/\x});

    % 最小値の点
    \draw[densely dotted, thick, gray!70!black] (5, 0) node[below, font=\footnotesize] {$5$}
    -- (5, 3.75)
    -- (0, 3.75) node[left, font=\footnotesize] {$\frac{15}{4}$};
    \draw[densely dotted, thick, gray!70!black] (6, 0) node[below, font=\footnotesize] {$6$}
    -- (6, 4.167)
    -- (0, 4.167) node[left, font=\footnotesize] {$\frac{25}{6}$};

    \node[circle, fill=red!85!black, inner sep=1.4pt] at (5, 3.75) {};
    \node[circle, fill=blue!85!black, inner sep=1.2pt] at (6, 4.167) {};

    % ラベル
    \node[right, font=\footnotesize, blue!85!black] at (3.5, 4.8) {$y=f(t)$};
  \end{tikzpicture}
  \caption{$y=f(t)$ の概形}
  \label{fig:graph_ft}
\end{figure}

従って\cref{eq:1}が$0<t<6$に解を持つ条件は
\begin{align}
  \dfrac{15}{4}\le u
\end{align}
であるから，求める最小値は$\dfrac{15}{4}\ \mathrm{m/s}$である．$\cdots$(答)

\end{document}
